\documentclass[11pt,onecolumn]{article}

% ── Packages ──────────────────────────────────────────────────────────────────
\usepackage[margin=1in]{geometry}
\usepackage{amsmath,amssymb}
\usepackage{graphicx}
\usepackage{booktabs}
\usepackage{hyperref}
\usepackage{natbib}
\usepackage{microtype}
\usepackage{xcolor}
\usepackage{lineno}
\usepackage{setspace}
\usepackage{caption}
\usepackage{subcaption}
\usepackage{authblk}
\usepackage{array}
\usepackage{tabularx}
\usepackage{orcidlink}

\hypersetup{
  colorlinks=true,
  linkcolor=blue!60!black,
  citecolor=blue!60!black,
  urlcolor=blue!60!black,
}

\doublespacing
\linenumbers

% ── Title block ───────────────────────────────────────────────────────────────
\title{\textbf{Mouse V1 Passes Through a Transient Low-Dimensional State
of Peak Cross-Individual Representational Consistency During Natural
Image Viewing}}

\author[1]{Z. Jana\orcidlink{0009-0009-3316-7476}}
\affil[1]{Department of Computer Science and Engineering,
The American University in Cairo}

\date{}

% ──────────────────────────────────────────────────────────────────────────────
\begin{document}
\maketitle

% ── Abstract ──────────────────────────────────────────────────────────────────
\begin{abstract}
When different mice view the same natural image, their primary visual
cortex (V1) population responses share a measurable fraction of
representational structure.  Using millisecond-resolution representational
similarity analysis (RSA) applied to Neuropixels recordings from all
32 \textit{brain\_observatory\_1.1} sessions in the Allen Brain
Observatory, we show that this cross-individual consistency is not
static: it follows a stereotyped two-component temporal trajectory,
rising to a transient peak at approximately 100~ms post-stimulus onset
(mean Spearman $r = 0.307$) before decaying to a lower sustained floor
($r = 0.182$, 220--250~ms window).  This transient peak coincides
precisely with a collapse in the effective dimensionality of the shared
representation: the participation ratio drops from a pre-stimulus
baseline of $\approx$35.6 to $\approx$15.3 at 100~ms within individual
mice ($\approx$108 to $\approx$31.5 in the cross-mouse-averaged
representation). Dimensionality and
consistency are tightly coupled across the entire post-stimulus epoch
(Spearman $r = -0.853$ across overlapping windows; corroborated by a
non-overlapping-window control, $r = -0.929$, $p = 4.6 \times 10^{-6}$):
the moment mice agree most is the moment the population code is most
compressed.  This
transient low-dimensional state is specific to complex, naturalistic
stimuli --- static gratings produce a curve that rises and then stays
essentially flat, with no transient peak and a near-unity
transient-to-sustained (T/S) ratio, while a control that matches natural
scenes to the gratings' item count confirms this shape difference is not
an artefact of the two stimulus sets having different numbers of
conditions --- and is highest in V1 itself (peak $r = 0.307$), with lower and
roughly comparable peak consistency across four extrastriate areas
(VISl $r = 0.186$, VISal $r = 0.234$, VISpm $r = 0.210$,
VISam $r = 0.238$).

To explain this transient compression, we tested whether it is
preferentially produced by fast-spiking inhibitory interneurons.
Waveform-based dissociation of putative fast-spiking (FSU,
narrow-waveform) and regular-spiking (RSU, broad-waveform) units ---
classified via a two-component Gaussian mixture fit to log waveform
duration, with a symmetric probability-mass bound excluding both the
ambiguous gap between the two modes and each class's own long tail,
rather than a single hand-picked cutoff --- shows that FSU units produce
a sharper transient than RSU units (T/S ratio 2.29$\times$ vs.\
1.71$\times$), a difference that survives subsampling RSU sessions down
to the FSU unit count (matched T/S $= 2.22\times$).  Sessions with
higher FSU fractions show correspondingly sharper population-level
transients ($\Delta$T/S $= +0.32$), and matching the low-FSU group's
unit count down to the high-FSU group's shows that most, but not all, of
the raw absolute-consistency gap between the groups was a similarity-matrix
SNR artifact of unit count rather than genuine FSU-driven signal.
Optotagged PV$^+$ interneurons from Pvalb-IRES-Cre$\times$Ai32 sessions
show an even sharper transient (T/S $= 3.59\times$) than waveform-inferred
FSU units (T/S $= 2.29\times$), consistent with partial impurity of the
waveform classification (of 18 optotagged PV$^+$ units, 6 fall in the
confidently-FSU range, 10 are excluded as ambiguous, and 2 fall in the
confidently-RSU range).
A Wilson-Cowan excitatory-inhibitory (E/I) firing-rate model fit directly
to the image-selective firing-rate amplitude of RSU and FSU populations
recovers an inhibition-dominated connectivity ratio
($r_\mathrm{EI} = w_\mathrm{EE}/w_\mathrm{IE} = 0.646$;
$R^2_\mathrm{RSU} = 0.901$, $R^2_\mathrm{FSU} = 0.900$) that generalises
to held-out sessions in five-fold cross-validation (mean test
$R^2 = 0.877 \pm 0.042$; $r_\mathrm{EI} < 1$ in every fold; mean
$r_\mathrm{EI} = 0.647 \pm 0.010$).  An independent fit of the same
model to cross-mouse consistency curves --- a different quantity
entirely --- converges on a similar ratio ($r_\mathrm{EI} = 0.608$),
whereas a fit to uncorrected total firing rate, which retains
non-selective drive, yields a markedly different, less stable, and
worse-fitting estimate ($r_\mathrm{EI} = 0.924$,
$R^2_\mathrm{FSU} = 0.601$), indicating the agreement between the first
two fits is not an artefact of curve-fitting flexibility.  Together,
these results show that mouse V1 passes through a transient,
inhibition-dominated low-dimensional state during natural image viewing,
and that this state is precisely when representational structure is
most shared across individuals.
\end{abstract}

\noindent\textbf{Keywords:} primary visual cortex; representational
similarity analysis; cross-individual variability; fast-spiking
interneurons; Wilson-Cowan model; neural geometry; Allen Brain
Observatory

\newpage

% ── Introduction ──────────────────────────────────────────────────────────────
\section{Introduction}

A central goal of systems neuroscience is to identify which features of
neural population activity are canonical --- shared across individuals of
the same species exposed to the same stimuli --- and which are
idiosyncratic, reflecting individual developmental history, anatomy, or
ongoing internal state.  In primary visual cortex, population responses
to natural images vary substantially across individuals in terms of
absolute firing rates and the identity of which neurons are active, yet
downstream circuits must extract a stable visual representation
regardless.  One quantitative handle on this problem is cross-individual
representational consistency: the degree to which different animals agree
on the relational structure of their population responses --- which
stimuli are represented as similar to which others --- even when their
absolute responses differ \citep{kriegeskorte2008}.  High cross-individual
consistency implies that a canonical representational geometry is expressed
in V1, potentially reflecting shared tuning properties that emerge from
common developmental programs and anatomical connectivity.

The geometry of V1 population representations has been characterised in
terms of dimensionality \citep{stringer2019}, the alignment of neural
manifolds across animals \citep{williams2021}, and the degree to which
response structure is predicted by deep convolutional neural network (CNN)
features \citep{yamins2016,cadena2019}.  These studies have generally
treated representational consistency as a static quantity, computed over a
broad temporal window that collapses the entire visual response into a
single point estimate.  However, visual responses in V1 are dynamically
structured: an early feedforward volley arriving approximately 40--70~ms
post-onset is followed by recurrent and feedback contributions that unfold
over the subsequent 200~ms \citep{lamme2000,girard2001}.  Whether
cross-individual agreement on representational structure is uniformly
distributed across this epoch, or concentrated in a particular temporal
component reflecting a specific circuit mechanism, is unknown.

Here we ask whether the V1 population code passes through a transient,
low-dimensional regime during the visual response, and if so, whether
that is the moment at which different animals' representations are most
alike.  A code transiently compressed into a small number of stereotyped
dimensions is, almost by construction, more likely to look similar across
individuals than a code spread across many idiosyncratic dimensions.
Cross-individual RSA consistency is therefore not only a measure of
agreement in its own right, but a sensitive, model-free readout of when
the population code is in such a compressed, canonical state.

Having established whether and when such a state occurs, we then ask
what produces it.  Fast-spiking interneurons --- the dominant subtype
of which expresses parvalbumin (PV$^+$) --- are a strong mechanistic
candidate.  PV$^+$ cells receive strong thalamic drive, fire with short
latency and high spike-timing precision \citep{swadlow2003}, and provide
powerful lateral inhibition onto pyramidal cells through perisomatic
synapses \citep{freund1996}.  The classical iceberg effect posits that
inhibition sculpts a broad excitatory input into a sharper, more
selective response \citep{somers1995,troyer1998}.  If this inhibitory
sharpening operates at the population level, fast-spiking units should
show a sharper version of any transient compression than broad-spiking
units, and a circuit model in which inhibition dominates recurrent
excitation should reproduce the observed dynamics.

We test these two questions in sequence using the Allen Brain Observatory
Visual Coding Neuropixels dataset \citep{siegle2021}, which provides
simultaneous recordings from hundreds of V1 units across 32 mice viewing
an identical set of 118 natural images.  We first apply a sliding-window
RSA approach at 10~ms temporal resolution to characterise the full
temporal trajectory of cross-individual consistency, and pair it with
participation ratio analysis to test directly whether peak consistency
coincides with minimum dimensionality.  We then turn to mechanism:
waveform-based cell-class dissociation, optotagged PV$^+$ recordings,
and a Wilson-Cowan E/I firing-rate model fit directly to firing rates,
cross-checked against an independent fit to consistency curves, and
cross-validated on held-out mice.  As mechanistic contrasts, we compare
the temporal
profiles of natural scenes with static gratings and examine how
consistency varies across five visual areas.

% ── Results ───────────────────────────────────────────────────────────────────
\section{Results}

\subsection{Cross-mouse consistency in V1 is substantial, varies across
natural images, and is predicted by response sparsity}

We recorded spike times from VISp units in 32 mice viewing 118 natural
images (250~ms each, 50 repetitions per image) from the Allen Brain
Observatory \textit{brain\_observatory\_1.1} dataset \citep{siegle2021}.
For each session, units passing quality control (amplitude cutoff $< 0.1$,
presence ratio $> 0.9$, ISI violation rate $< 0.5$; $n = 14$--$110$ units
per session after filtering) were used to construct a
$118 \times n_\mathrm{units}$ trial-averaged, column-wise z-scored response
matrix.  Cross-mouse consistency was quantified using representational
similarity analysis (RSA): for each image, we computed the cosine
similarity of its population response vector to all other 117 images
within each mouse, then measured how much mice agreed on this ranking
using the mean Spearman rank correlation across all
$\binom{32}{2} = 496$ mouse pairs.

Within-mouse split-half reliability of the raw response vectors
(Spearman-Brown corrected, 20 random splits) was near unity
(mean $r = 0.978$, range $[0.968, 0.983]$; Figure~\ref{fig:consistency}A),
confirming that the analysis is not power-limited by within-session
trial noise.  This response-vector reliability is not directly comparable
to cross-mouse consistency, however, since the two are computed in
different spaces: reliability here is a per-unit response-vector
correlation, whereas consistency is computed on image~$\times$~image
similarity structure (RDM space).  Repeating the split-half procedure in
RDM space --- building two similarity matrices from independent trial
halves and correlating them exactly as cross-mouse consistency
correlates two mice's similarity matrices --- gave a within-mouse
RDM-space reliability of mean $r = 0.832$ (range $[0.609, 0.929]$),
which is the number properly comparable to the consistency values
reported below.

\begin{figure}[htbp]
  \centering
  \includegraphics[width=\textwidth]{figures/fig1_consistency_overview.png}
  \caption{\textbf{Cross-mouse consistency is substantial but well below
  within-mouse reliability.}
  \textbf{(A)} Distribution of Spearman-Brown corrected split-half
  response-vector reliability across 118 natural images, averaged across
  32 mice (mean $= 0.978$, range $[0.968, 0.983]$).  Near-ceiling values
  confirm that cross-mouse consistency is not limited by within-session
  trial noise.  Dashed line: mean.
  \textbf{(B)} Mean Spearman $r$ across 496 mouse pairs for each of the
  118 natural images, sorted from highest to lowest.  Shading: 95\%
  cluster bootstrap confidence intervals (resampling whole mice, 1000
  iterations; self-pairs excluded).  Mean $r = 0.314$; range
  0.142--0.588.  Dashed horizontal line: within-mouse RDM-space
  reliability (mean $r = 0.832$), the reliability estimate computed in
  the same space as consistency.  All images are significantly above
  zero.}
  \label{fig:consistency}
\end{figure}

Cross-mouse consistency was significantly above zero for all 118 images
(mean $r = 0.314$, range 0.142--0.588;
Figure~\ref{fig:consistency}B).  The fourfold range across images
indicates that some images drive strongly shared representations while
others are encoded more idiosyncratically.

Response sparsity --- the fraction of units responding below the
per-mouse median absolute z-score --- was an image-level predictor of
consistency (Spearman $r = -0.470$,
$p = 7.8 \times 10^{-8}$): images activating fewer units were more
consistently represented across mice.  Response magnitude was a
secondary predictor (Spearman $r = +0.276$, $p = 2.5 \times 10^{-3}$;
Figure~\ref{fig:predictors}A--B).  A ResNet50 CNN feature embedding
also predicted consistency in leave-one-out ridge regression; comparing
across layers, layer~2 (512-dimensional) gave the best accuracy
(LOO Pearson $r = -0.496$, $p = 1.1 \times 10^{-8}$) and was selected as
the primary CNN predictor, with the final layer (1000-dimensional)
performing comparably ($r = -0.486$, $p = 2.5 \times 10^{-8}$;
Figure~\ref{fig:predictors}C).  To ask whether this CNN-based signal was
subsumed by sparsity, we compared nested regression models: adding the
CNN prediction to a sparsity-only model ($R^2 = 0.269$) raised $R^2$ to
0.453 for layer~2 ($\Delta R^2 = +0.184$) and to 0.374 for the final
layer ($\Delta R^2 = +0.105$; Figure~\ref{fig:predictors}D).  This
nested-$R^2$ comparison is the primary evidence that CNN features carry
signal beyond sparsity, since it is a direct model-fit comparison rather
than a residual-alignment statistic.  As a secondary check, added-variable
partial correlation, controlling for sparsity, retained
$r_\mathrm{partial} = -0.501$ for layer~2 and $-0.378$ for the final
layer.  A second, non-parametric partialling method (rank-residualised
Spearman correlation) gave a substantially larger-magnitude partial
correlation for both layers ($r_\mathrm{partial} = -0.828$,
$p = 5.5\times10^{-31}$ for layer~2; $r_\mathrm{partial} = -0.887$,
$p = 9.9\times10^{-41}$ for the final layer; Supplementary
Figure~S2C).  This magnitude jump across partialling methods is a
signature of sparsity acting as a partial suppressor --- it shares
variance with both consistency and the CNN prediction, so removing it
from both can inflate residual alignment --- and is a reason for caution
rather than confidence: we treat the partial-correlation values as
illustrative rather than as precise effect sizes, and rely on the
nested $\Delta R^2$ for the headline claim (Methods).  Response
sparsity and CNN-layer features therefore each explain distinct
variance in cross-mouse consistency, and this conclusion is robust to
the specific CNN layer used.  Galleries of the most and least
consistently represented images confirm the sparsity relationship
qualitatively: high-consistency images tend to be sparse, high-contrast,
or dominated by a single feature, while low-consistency images are
denser and more spatially heterogeneous (Supplementary Figure~S1).

\begin{figure}[htbp]
  \centering
  \includegraphics[width=\textwidth]{figures/fig2_image_predictors.png}
  \caption{\textbf{Response sparsity and CNN-layer features each explain
  distinct variance in cross-mouse consistency.}
  \textbf{(A)} Consistency vs.\ response sparsity (fraction of units
  below per-mouse median $|z|$; Spearman $r = -0.470$,
  $p = 7.8 \times 10^{-8}$).  Line: linear reference fit.
  \textbf{(B)} Consistency vs.\ mean $|$z-scored response$|$
  (Spearman $r = +0.276$, $p = 2.5 \times 10^{-3}$).
  \textbf{(C)} LOO Pearson $r$ per ResNet50 layer (purple: layer~2,
  the primary CNN predictor; teal: the final layer, the secondary
  predictor; grey: other layers).  The final
  layer performs comparably ($r = -0.486$) and is carried through
  panel~D and Supplementary Figure~S2 as a robustness check.
  \textbf{(D)} Nested $R^2$: a sparsity-only model ($R^2 = 0.269$)
  compared against sparsity~$+$~layer~2 ($R^2 = 0.453$,
  $\Delta R^2 = +0.184$) and sparsity~$+$~final layer ($R^2 = 0.374$,
  $\Delta R^2 = +0.105$).  This model-comparison statistic, rather than
  the partial-correlation coefficient, is the primary evidence that CNN
  features add explanatory power beyond sparsity; partial-correlation
  estimates, which vary substantially with partialling method, are
  reported as a robustness check in Supplementary Figure~S2.}
  \label{fig:predictors}
\end{figure}

\subsection{Cross-mouse consistency follows a stereotyped two-component
temporal trajectory}

To characterise how cross-mouse agreement evolves over the visual
response, we applied the same RSA pipeline to 25~ms sliding windows
stepped at 10~ms increments from $-50$ to $+370$~ms relative to
stimulus onset.  Pre-stimulus consistency was near zero
(mean $r = 4 \times 10^{-5}$), confirming that the measure is not
contaminated by correlated spontaneous activity or slow drift across
sessions (Figure~\ref{fig:temporal}).

Following stimulus onset, consistency rose steeply from $\sim$40~ms,
reaching its first threshold-crossing value at 50~ms and peaking at
$r = 0.307$ at 100~ms.  Consistency then decayed monotonically to a
sustained floor estimated from the 220--250~ms window
($r = 0.182 \pm 0.003$; one-sample $t$-test vs.\ zero:
$t = 94.75$, $p = 2.6 \times 10^{-6}$).  We restrict the sustained
floor estimate to the 220--250~ms epoch specifically to avoid
contamination from the next image's neural response: in the
\textit{brain\_observatory\_1.1} protocol, natural scenes are presented
back-to-back with no inter-stimulus interval, so the subsequent image's
V1 response arrives at approximately 290--300~ms (stimulus onset at
250~ms plus $\approx$40--50~ms response latency).  Data beyond 250~ms
are not interpreted as reflecting the sustained representation of the
current image.

A coarser
comparison of early (0--75~ms) and late (75--250~ms) windows confirmed
that late-window consistency substantially exceeds early-window
consistency (mean $r = 0.084$ vs.\ $0.320$; Wilcoxon signed-rank
$p = 4.2 \times 10^{-21}$; Supplementary Figure~S3), while image rankings
were partially preserved across epochs (Spearman $r = 0.652$ between
the 100~ms and 250~ms per-image consistency vectors; Supplementary
Figure~S7).

\begin{figure}[htbp]
  \centering
  \includegraphics[width=\textwidth]{figures/fig3_temporal_profile.png}
  \caption{\textbf{Cross-mouse consistency follows a stereotyped
  two-component trajectory.}
  Mean cross-mouse consistency over time (25~ms windows, 10~ms steps,
  $n = 32$ mice; blue).  Pre-stimulus baseline $\approx 4 \times 10^{-5}$,
  validating the measure.  Onset: 50~ms; peak: $r = 0.307$ at 100~ms
  (green marker).  Sustained floor: $r = 0.182$ (220--250~ms, shaded
  band and dashed horizontal line).  Vertical dashed line: stimulus onset
  (0~ms); dotted line: stimulus offset (250~ms), beyond which
  back-to-back presentation creates contamination risk from the next
  image's V1 response ($\sim$290--300~ms onset).}
  \label{fig:temporal}
\end{figure}

\subsection{The transient peak coincides with a collapse of the shared
representation into a low-dimensional attractor}

To characterise the geometry of the shared representation at each time
point, we computed the participation ratio (PR) of the average cosine
similarity matrix across mice:
\begin{equation}
  \mathrm{PR} = \frac{\bigl(\sum_i \lambda_i\bigr)^2}{\sum_i \lambda_i^2},
  \label{eq:PR}
\end{equation}
where $\lambda_i$ are the non-negative eigenvalues of the average
similarity matrix.  PR measures effective dimensionality: PR~$= 1$ if one
eigenvalue dominates; PR~$= n$ if all eigenvalues are equal.

Pre-stimulus PR of the cross-mouse-averaged similarity matrix was
$\approx 108$ --- close to the maximum of 118 given the number of
images --- indicating that spontaneous activity produces an
approximately uniform representational geometry with no shared structure
across mice.  Because averaging similarity matrices across mice pushes
disagreement toward uniformity, this averaged-matrix PR is partly
definitional with cross-mouse consistency; a less confounded view is
the PR of each mouse's own similarity matrix, averaged after the fact
(pre-stimulus per-mouse PR $\approx 35.6$), which never involves
cross-mouse averaging.  At the consistency peak (100~ms window), the
averaged-matrix PR dropped to 31.5, meaning the shared representation
occupied only $\approx$27\% of the available dimensions; the per-mouse
PR dropped correspondingly, from $\approx 35.6$ pre-stimulus to 15.3 at
the peak.  PR then partially recovered to 61.2 (averaged-matrix) and
24.4 (per-mouse) at 250~ms, tracking the decay of consistency
(Figure~\ref{fig:geometry}A).  Across all post-stimulus windows, PR and
mean consistency were strongly negatively correlated both for the
averaged-matrix PR (Spearman $r = -0.853$, $p = 3.5 \times 10^{-6}$) and,
essentially as strongly, for the less-confounded per-mouse PR
(Spearman $r = -0.805$, $p = 3.2 \times 10^{-5}$): the moment of maximal
cross-mouse agreement is precisely the moment of maximal representational
compression, and this holds even when the compression is measured
entirely within individual mice.

Because these post-stimulus windows overlap by up to 60\% of their width
(25~ms windows stepped by 10~ms), the $n \approx 19$ points entering this
correlation are not independent, and the analytic $p$-values above
understate the true $p$-value (as similarly noted for Supplementary
Figure~S7B). We assessed this two ways (Methods). First, a circular-shift
permutation test, which preserves each series' own autocorrelation while
destroying only their mutual alignment, gave $p = 0.053$ for both the
averaged-matrix and per-mouse correlations --- the observed $r$ was more
extreme than every one of the 18 possible shifts, so this is the floor
$p$-value achievable with this few independent windows, not evidence of a
weak effect. Second, we recomputed the correlation on windows spaced
30~ms apart (one window-width, so no window shares spike-count bins with
its neighbour): this fully independent, 13-point version gave
$r = -0.929$, $p = 4.6 \times 10^{-6}$, confirming the coupling under a
conventional test that needs no autocorrelation correction.

Because PR and consistency are both computed from the same response
matrices, and both are plausibly driven by the stimulus-evoked response
envelope (weak pre-stimulus and late, strong near the transient peak), we
tested whether this coupling is better explained as a shared dependence on
overall response magnitude (Methods). The
envelope --- mean raw population response per window, independent of the
per-window z-scoring used for RSA --- itself correlated with both
consistency ($r = 0.525$, analytic $p = 0.021$, circular-shift
$p = 0.211$) and PR (avg-matrix: $r = -0.711$, analytic
$p = 6.5\times10^{-4}$, circular-shift $p = 0.053$; per-mouse:
$r = -0.653$, analytic $p = 2.5\times10^{-3}$, circular-shift $p = 0.053$).
Under the autocorrelation-robust test, the envelope--consistency
correlation is no longer distinguishable from chance, while the
PR--envelope correlations remain at the floor $p$-value achievable given
$n \approx 19$ independent windows (Methods); if anything, this weakens
rather than strengthens a common-cause account of the PR--consistency
coupling. The PR--consistency coupling itself survived partialling out
the envelope nearly unchanged (avg-matrix: partial $r = -0.801$, analytic
$p = 6.5\times10^{-5}$, circular-shift $p = 0.053$; per-mouse: partial
$r = -0.718$, analytic $p = 8.0\times10^{-4}$, circular-shift $p = 0.053$,
vs.\ raw $r = -0.805$), indicating that PR tracks consistency well beyond
what the response envelope predicts on its own.

Eigendecomposition of the average similarity matrix at 100~ms revealed
that the top 8 eigenvectors account for 35.2\% of variance, with a
gradually decaying spectrum (14.75\%, 4.60\%, 3.57\%, 3.21\%, 2.77\%,
2.12\%, 2.07\%, 2.07\% for dimensions 1--8; Figure~\ref{fig:geometry}B).
The two dominant dimensions were interpretable in terms of low-level
image statistics: Dim~1 correlated positively with response magnitude
(Spearman $r = +0.44$) and early CNN feature strength
(layer~1: $r = +0.41$; layer~4: $r = -0.36$), representing a
signal-strength axis; Dim~2 correlated with response sparsity
($r = -0.51$), representing a spatial-scale axis (see Supplementary
Figure~S5 for the full property heatmap).  Consistent with the
distributed nature of the shared structure, overall cross-mouse
consistency did not strongly load onto any single eigenvector
(maximum $|r| = 0.28$; Figure~\ref{fig:geometry}C).  A second-order RSA
--- correlating each mouse's full representational dynamics profile
(similarity to all other images across five time windows) between mouse
pairs --- recovered a mean second-order $r = 0.187$ across images and a
temporal peak at exactly 100~ms.  This is a related but distinct
consistency statistic rather than an independent confirmation: it
correlates with first-order consistency at $r = 0.950$ across images
(Supplementary Figure~S8C), so it is better read as a re-parameterisation
that tracks the same underlying signal than as an orthogonal check.

\begin{figure}[htbp]
  \centering
  \includegraphics[width=\textwidth]{figures/fig4_representational_geometry.png}
  \caption{\textbf{The consistency peak coincides with collapse into a
  low-dimensional attractor.}
  \textbf{(A)} Participation ratio (PR; right axis, light purple) and mean
  cross-mouse consistency (left axis, blue) co-plotted over time.
  Averaged-matrix PR: pre-stimulus $\approx 108$; peak-window PR $= 31.5$
  (27\% of available dimensions); 250~ms PR $= 61.2$.  Per-mouse PR
  (shaded band, $\pm$ SEM): pre-stimulus $\approx 35.6$; peak-window
  PR $= 15.3$; 250~ms PR $= 24.4$.  Post-stimulus Spearman
  $r(\text{PR, consistency})$ annotated on panel for both the
  averaged-matrix PR ($-0.853$, $p = 3.5 \times 10^{-6}$) and the
  per-mouse PR ($-0.805$, $p = 3.2 \times 10^{-5}$).
  \textbf{(B)} Variance explained by the top 8 eigenvectors of the
  average cosine similarity matrix at 100~ms (bars, left axis;
  cumulative line, right axis; top 8 dims $= 35.2\%$ total).
  \textbf{(C)} Image layout in two-dimensional attractor space
  (Dim~1: signal-strength axis; Dim~2: sparsity axis), coloured by
  cross-mouse consistency.  Consistency does not concentrate at the
  extremes of either axis, indicating the shared structure is distributed
  throughout the image set.}
  \label{fig:geometry}
\end{figure}

\subsection{Fast-spiking interneurons drive the transient component of
the shared attractor}

To test whether the transient and sustained components of cross-mouse
consistency are differentially produced by distinct cell classes, we
classified VISp units by waveform trough-to-peak duration.  Rather than a
single hand-picked cutoff, we fit a 2-component Gaussian mixture model
(GMM) to log waveform duration across all 3,694 QC-passing VISp units
(fitting in log-space prevents a small tail of long-duration units,
$\sim$1--2.2~ms, likely triphasic/artifact waveforms, from being fit as
its own spurious component instead of the true narrow/broad structure).
The fit recovered component means of 0.263~ms (narrow) and 0.627~ms
(broad).  A unit was classified as putative fast-spiking (FSU) or
regular-spiking (RSU) only if (a) its posterior probability of component
membership exceeded 0.9, excluding the ambiguous gap between the two
modes, and (b) it fell within a symmetric interval capturing 60\% of that
component's own probability mass around its mean (0.84 standard
deviations in log-space), excluding the long tail from being swept into
RSU merely for being broader than FSU.  This yielded an FSU range of
0.218--0.316~ms and an RSU range of 0.549--0.728~ms (mean 8.0 and 39.4
units per session, respectively; 28.2\% of units fell in the excluded
ambiguous/outlier band).  We then computed cross-mouse consistency
separately for each subpopulation at each sliding-window time point
(minimum 8 units per session per class; FSU: 15/32 sessions; RSU: 31/32
sessions).

\subsubsection{Waveform-based dissociation}

RSU and FSU subpopulations rose with identical onset latency ($\sim$40~ms),
ruling out a temporal-precedence account of the cell-class difference
(Figure~\ref{fig:celltypes}A).  FSU units showed a higher transient peak
($r = 0.245$) than RSU units ($r = 0.217$) and a lower sustained floor
($r = 0.107$ vs.\ $0.127$), yielding a substantially sharper T/S ratio:
$2.29\times$ for FSU vs.\ $1.71\times$ for RSU (Figure~\ref{fig:celltypes}B).
When all units were combined, cross-mouse consistency exceeded either
subpopulation alone at the peak, suggesting superadditive contributions
from the two cell classes encoding complementary image dimensions.

\subsubsection{Unit-count-matched control}

RSU sessions have roughly five times as many units as FSU sessions
(mean 39.4 vs.\ 8.0), which inflates similarity-matrix signal-to-noise
ratio, and hence raw consistency, independently of any genuine cell-type
effect.  To isolate the cell-type contribution from this confound, we
randomly subsampled each RSU session's unit pool down to 8 units
(matching the mean FSU count) over 50 draws and recomputed the sliding
consistency curve and T/S ratio on each matched draw.  The unit-matched
RSU T/S ratio ($2.22\times$; draws: mean $2.23\times$, 95\% range
$1.99$--$2.56\times$) remained below the FSU T/S ratio ($2.29\times$),
confirming that the FSU/RSU T/S difference is not simply an artifact of
RSU sessions having more units, though the matched control also narrows
the gap relative to the raw comparison ($1.71\times$), indicating unit
count does contribute part of the raw effect size.

\subsubsection{Session-level FSU fraction}

To test whether between-session variation in FSU abundance predicted the
sharpness of the population-level transient, we split sessions at the
median FSU fraction into high (mean fraction 0.255, $n = 16$) and low
(mean fraction 0.135, $n = 16$) groups (Figure~\ref{fig:celltypes}C).
The high-FSU group showed a T/S ratio of $1.88\times$ vs.\ $1.57\times$
for the low-FSU group (difference $+0.32$), with the high-FSU group
peaking more sharply and decaying faster, consistent with FSU units being
the primary driver of the transient component.

Because low-FSU sessions mechanically have more total (mostly RSU) units
--- FSU count is roughly stable while the denominator shrinks --- the
low-FSU group also has a similarity-matrix SNR advantage, which produced
a higher raw absolute consistency for the low-FSU group (mean $0.212$ vs.\
$0.155$ for high-FSU over the 0--250~ms window, a gap of $+0.057$) that is
not itself evidence against the FSU-driven transient account.  Subsampling
the low-FSU group's unit pool down to the high-FSU group's mean unit
count (55/session) reduced this gap to $+0.043$ ($0.198$ vs.\ $0.155$),
showing that roughly a quarter of the raw absolute-consistency gap is an
SNR artifact of unit count rather than a genuine FSU effect, while a
residual gap remains; the T/S ordering (high-FSU $> $ low-FSU) held on the
matched groups as well ($1.88\times$ vs.\ $1.62\times$).

\subsubsection{\texorpdfstring{Optotagged PV$^+$ validation and the limits of waveform
classification}{Optotagged PV+ validation and the limits of waveform classification}}

To directly validate the waveform-based FSU classification and characterise
the T/S behaviour of genetically confirmed PV$^+$ units, we identified
optotagged interneurons from the five Pvalb-IRES-Cre/wt$\times$Ai32
sessions in the \textit{brain\_observatory\_1.1} dataset.  Units were
classified as optotagged PV$^+$ if they met two criteria simultaneously:
mean spike count in the 0--5~ms post-pulse window $> 0.05$ spikes per
pulse, and light-to-baseline ratio $> 2.0$ (where baseline is the 5~ms
pre-pulse window).  The 5~ms response window was chosen to exclude
polysynaptic responses.

This procedure yielded 18 optotagged VISp units across five sessions.
Three sessions contributed $\geq 3$ optotagged units (session IDs:
721123822, 760345702, 797828357; contributing 3, 3, and 10 confirmed
PV$^+$ units respectively) and were included in the cross-mouse RSA;
two sessions contributed only one unit each and were excluded.

A critical complication emerged from the waveform durations of the
optotagged units, and the two-threshold GMM classification (above)
resolves it more informatively than a single cutoff would.  Under the
GMM's confidently-classified ranges (FSU: 0.218--0.316~ms; RSU:
0.549--0.728~ms), only 6 of the 18 optotagged units fall in the
confidently-FSU range, 2 fall in the confidently-RSU range (durations
0.577 and 0.618~ms, both from session 760345702), and the remaining 10
fall in the excluded ambiguous/outlier band --- i.e., most optotagged
PV$^+$ units have waveform durations too close to the narrow/broad
boundary to be confidently classified either way, rather than being
cleanly but wrongly assigned to RSU.  This is a more accurate picture
than the single-threshold analysis it replaces (which forced 15 of 18
into FSU and 3 into RSU) and indicates that the waveform-inferred FSU
category, even under a principled two-threshold rule, is enriched for
but not identical to the confirmed PV$^+$ population.  Quantitatively,
this residual contamination deflates the FSU T/S ratio slightly (by
excluding some genuine PV$^+$ cells to the ambiguous band, or in two
cases assigning them to RSU) relative to the true PV$^+$ T/S ratio.

This prediction is borne out: confirmed PV$^+$ units show a sharper
transient (T/S $= 3.59\times$) than waveform-inferred FSU units
(T/S $= 2.29\times$; Figure~\ref{fig:celltypes}D--E).  The direction
and magnitude of the effect are conserved and strengthened when the
cell-type attribution is made by genetics rather than waveform, even with
only three sessions contributing to the optotagged analysis.  The absolute
consistency values are lower in the optotagged analysis (fewer units per
session $\to$ noisier similarity matrices), but T/S is a shape property
of the normalised curve and is unaffected by overall signal level.

\begin{figure}[htbp]
  \centering
  \includegraphics[width=\textwidth]{figures/fig5_cell_type_dissociation.png}
  \caption{\textbf{Fast-spiking, genetically confirmed PV$^+$ interneurons
  produce a sharper transient than broad-spiking units.}
  \textbf{(A)} Temporal consistency curves for RSU (pink), FSU (teal),
  and all units (dashed grey).  RSU and FSU rise with identical onset
  latency ($\sim$40~ms; annotated arrow) but diverge in T/S ratio
  (FSU: 2.29$\times$; RSU: 1.71$\times$).  Triangle markers: peak;
  dashed horizontal lines: sustained floor for each class.
  \textbf{(B)} T/S ratio bar chart per cell class.  Paired bars: peak
  (opaque) and sustained floor (semi-transparent); T/S ratio annotated
  above each pair.
  \textbf{(C)} Sessions split at median FSU fraction.  High FSU
  ($n = 16$, mean fraction $= 0.255$, T/S $= 1.88\times$, teal-green) vs.\
  low FSU ($n = 16$, mean fraction $= 0.135$, T/S $= 1.57\times$, dark red;
  $\Delta$T/S $= +0.32$).  Dashed grey: all sessions.
  \textbf{(D)} Temporal curves for waveform-inferred FSU ($n = 15$
  sessions, teal) and optotagged PV$^+$ units (light blue dashed,
  $n = 3$ sessions, 18 confirmed units).  Lower absolute consistency in
  the optotagged analysis reflects fewer units per session; T/S is a
  shape property unaffected by signal level.
  \textbf{(E)} T/S ratio comparison: optotagged PV$^+$ $= 3.59\times$
  vs.\ waveform FSU $= 2.29\times$, consistent with partial impurity of
  the waveform classification (of 18 optotagged PV$^+$ units, 6 fall in
  the confidently-FSU range, 10 are excluded as ambiguous, and 2 fall in
  the confidently-RSU range).
  \textbf{(F)} Waveform duration distribution across all 3,694
  QC-passing VISp units, colored by classification: FSU (teal), RSU
  (pink), and excluded ambiguous/outlier units (grey).  Classes were
  assigned by a 2-component Gaussian mixture fit to log waveform duration
  with a symmetric per-component probability-mass bound (FSU range:
  0.218--0.316~ms; RSU range: 0.549--0.728~ms; 28.2\% of units excluded).}
  \label{fig:celltypes}
\end{figure}

\subsection{A Wilson-Cowan E/I model fit to firing rates converges with
an independent fit to consistency curves on inhibition-dominated
connectivity}

Every fit described so far treats cross-mouse consistency as the target
of interest, but consistency is an agreement statistic, not a firing
rate. A rate model's success in reproducing consistency curves does not,
on its own, establish that a rate model's E/I dynamics track anything
about actual population activity. To test this directly, we fit a
two-population Wilson-Cowan firing-rate model \citep{wilson1972} to
three different targets in turn: the image-selective firing-rate
amplitude of the RSU (E population) and FSU (I population) pools
(computed directly from the raw spike-count tensors), the same
populations' uncorrected total firing rate, and the RSU/FSU temporal
consistency curves used throughout the paper. Both populations receive a
shared gamma-shaped thalamic drive (onset $\sim$40~ms; shape fixed to
match the observed response latency).

The model is described by:
\begin{align}
  \tau_E \dot{E} &= -E + \sigma\!\left(r_\mathrm{EI}\, w_\mathrm{IE}\, E
    - w_\mathrm{IE}\, I + D(t)\right), \label{eq:E} \\
  \tau_I \dot{I} &= -I + \sigma\!\left(w_\mathrm{EI}\, E
    - w_\mathrm{II}\, I + 0.7\, D(t)\right), \label{eq:I}
\end{align}
where $\sigma(x) = 1/(1 + e^{-x})$ is the sigmoid activation function,
$D(t)$ is a gamma-shaped thalamic drive, $\tau_E = 20$~ms and
$\tau_I = 10$~ms are fixed at biologically motivated values, and the
five free parameters are
$\{r_\mathrm{EI}, w_\mathrm{IE}, w_\mathrm{EI}, w_\mathrm{II},
I_\mathrm{amp}\}$.  The key design choice is the explicit
reparameterisation of $r_\mathrm{EI} = w_\mathrm{EE}/w_\mathrm{IE}$ as
a free parameter (bounded $[0.01, 2.0]$): $r_\mathrm{EI} < 1$ indicates
inhibition-dominated dynamics; $r_\mathrm{EI} > 1$ indicates
recurrence-dominated dynamics.  The absolute scale of $w_\mathrm{IE}$
is a nuisance parameter not well-determined by the data; $r_\mathrm{EI}$
is the primary inference target.

Total firing rate mixes image-selective, stimulus-driven modulation with
non-selective drive shared identically across all images (e.g.\ a broad
contrast/onset transient). The latter is exactly what per-unit z-scoring
removes before consistency is computed (Methods), so total rate and
consistency are not measuring the same thing even in principle. We
therefore also computed an image-selective firing-rate amplitude for
each pool: the per-unit standard deviation of firing rate across images,
averaged over the pool, in the same sliding windows used for the
consistency curves. A non-selective burst identical across images
contributes $\approx 0$ to this quantity, mirroring the z-scoring step,
while remaining a rate-derived measurement rather than a cross-mouse
agreement statistic.

Fit directly to this image-selective amplitude, optimisation yielded
$r_\mathrm{EI} = 0.646$ ($R^2_\mathrm{RSU} = 0.901$,
$R^2_\mathrm{FSU} = 0.900$; Figure~\ref{fig:eimodel}A), with the fit
essentially uniquely identified across 10 random initialisations
(mean $r_\mathrm{EI} = 0.646$, SD $< 0.001$). Fit independently to the
cross-mouse consistency curves, the same model converged on a similar
ratio ($r_\mathrm{EI} = 0.608$, $R^2_\mathrm{RSU} = 0.901$,
$R^2_\mathrm{FSU} = 0.913$), though with more restart-to-restart
variability: across 10 random initialisations, 9 converged to the same
inhibition-dominated solution ($r_\mathrm{EI} \approx 0.608$) while one
escaped to a competing, recurrence-dominated basin near
$r_\mathrm{EI} \approx 1.01$ (mean $r_\mathrm{EI} = 0.648 \pm 0.120$
across all 10; Supplementary Figure~S9C). Because these two fits target
different quantities computed from largely non-overlapping information
--- a rate amplitude and a cross-mouse agreement statistic --- their
agreement is convergent evidence rather than a restatement of the same
result.

As a specificity check, fitting the same model to \textit{uncorrected}
total firing rate, which retains the non-selective drive component,
yielded a markedly different, less stable, and worse-fitting estimate
($r_\mathrm{EI} = 0.924$; $R^2_\mathrm{RSU} = 0.805$,
$R^2_\mathrm{FSU} = 0.601$; mean across restarts
$= 0.926 \pm 0.006$; Supplementary Figure~S9B, D). This confirms that
the agreement between the selective-amplitude and consistency fits is
not simply an artefact of a five-parameter model's flexibility to match
any pair of smooth curves: fitting the wrong signal produces a visibly
worse and quantitatively different answer. Sensitivity analysis
confirmed that T/S ratios $> 1$ require $r_\mathrm{EI} < 1$ within this
model family: increasing $r_\mathrm{EI}$ above 1 flattens the transient
and produces a sustained plateau inconsistent with the data
(Figure~\ref{fig:eimodel}B).

To test whether the inhibition-dominance conclusion generalises to
sessions not used in fitting, we performed five-fold cross-validation
over sessions on the selective-amplitude fit (Figure~\ref{fig:eimodel}C).
In each fold, RSU and FSU selective-amplitude curves were computed from
the $\sim$80\% training sessions and the model was fit using three warm
starts near the full-data solution; the fitted parameters were then used
to predict curves from the held-out $\sim$20\%.  Results: mean train
$R^2 = 0.899 \pm 0.006$; mean test $R^2 = 0.877 \pm 0.042$;
$r_\mathrm{EI} = 0.647 \pm 0.010$; $r_\mathrm{EI} < 1$ in every fold
(fold-level values: 0.635, 0.649, 0.660, 0.653, 0.640).  The
inhibition-dominance conclusion is robust even in the worst-performing
fold (test $R^2 = 0.804$, $r_\mathrm{EI} = 0.653$).  Because each fold is
fit with warm starts near the full-data solution, this cross-validation
demonstrates that the fit \textit{quality} generalises to held-out
sessions; it does not by itself
demonstrate that $r_\mathrm{EI}$ is basin-independent within each fold ---
that evidence comes instead from the 10-random-initialisation stability
analysis above.

\begin{figure}[htbp]
  \centering
  \includegraphics[width=\textwidth]{figures/fig6_ei_model.png}
  \caption{\textbf{A Wilson-Cowan model fit directly to image-selective
  firing rates recovers inhibition-dominated connectivity and
  generalises to held-out sessions.}
  \textbf{(A)} Model fit (dashed) vs.\ data (solid) for RSU and FSU
  image-selective firing-rate amplitude ($R^2_\mathrm{RSU} = 0.901$,
  $R^2_\mathrm{FSU} = 0.900$). $r_\mathrm{EI} = 0.646$.
  \textbf{(B)} T/S ratio sensitivity as a function of $r_\mathrm{EI}$
  (x-axis) and $w_\mathrm{IE}$ (y-axis), other parameters fixed at the
  selective-amplitude fit's values.  Star: fitted parameters.
  White dashed line: $r_\mathrm{EI} = 1$ (balance point between
  recurrent excitation and lateral inhibition).  High T/S ratios
  require $r_\mathrm{EI} < 1$ (left of balance line).
  \textbf{(C)} Five-fold cross-validation of the selective-amplitude fit
  over sessions.  Bars (left axis): train $R^2$ and test $R^2$ per fold;
  mean test $R^2 = 0.877 \pm 0.042$.  Points (right axis):
  $r_\mathrm{EI}$ per fold; mean $= 0.647 \pm 0.010$; reference at
  $r_\mathrm{EI} = 1$.  $r_\mathrm{EI} < 1$ in all five held-out subsets.
  \textbf{(D)} $r_\mathrm{EI}$ recovered from three independent fit
  targets --- cross-mouse consistency, uncorrected total firing rate,
  and image-selective firing-rate amplitude --- with error bars showing
  SD across 10 random restarts.  The selective-amplitude and consistency
  fits converge on similar values despite being computed from different
  quantities; the total-rate fit, which retains non-selective drive,
  gives a markedly different and less stable estimate.  See
  Supplementary Figure~S9 for the consistency and total-rate fit curves
  and the $w_\mathrm{II}$-bound sensitivity of all three fits.}
  \label{fig:eimodel}
\end{figure}

\subsection{Controls: adaptation and behavioural state do not explain
the results}

\paragraph{Adaptation.}
If all mice habituate to the same images in the same way, their responses
might converge over the session through shared adaptation rather than
shared tuning.  To test this, we split each session's trials by temporal
order and computed cross-mouse consistency from the chronologically first
and second halves separately.  Per-image consistency was highly stable
across halves (Spearman $r = 0.960$, $p = 6.1 \times 10^{-66}$;
Figure~\ref{fig:controls}A), and both halves correlated with full-session
consistency at $r > 0.984$ (Supplementary Figure~S10B--C).  The mean
per-image difference (early $-$ late consistency) was a small but
statistically significant $+0.015$ (one-sample $t$-test
$p = 1.1 \times 10^{-8}$) --- in the direction opposite to a shared
adaptation confound, which would predict convergence in the second half
--- consistent instead with modest non-selective neural fatigue that
attenuates all images equally without distorting representational
structure.

\paragraph{Behavioural state.}
If mice happen to be in more similar arousal or locomotion states during
certain images, those images would appear more consistently represented
for reasons unrelated to their visual content.  The coefficient of
variation (CV) of pupil area across mice per image was not significantly
correlated with per-image consistency (Spearman $r = +0.056$, $p = 0.55$,
$n = 26$ sessions with eye-tracking; Figure~\ref{fig:controls}B).  The CV
of running speed was likewise uncorrelated ($r = -0.033$, $p = 0.72$,
$n = 32$ sessions; Figure~\ref{fig:controls}C).  The narrow range of CV
across images (pupil CV SD $\approx 0.014$, running CV SD
$\approx 0.026$) limits statistical power for both tests, and
particularly for the pupil test given its smaller session count
($n = 26$); with this little variation to correlate against, the tests
are underpowered to detect anything but a fairly large true effect, so
the null result is interpreted cautiously rather than as positive
evidence of no relationship.

\begin{figure}[htbp]
  \centering
  \includegraphics[width=\textwidth]{figures/fig7_controls.png}
  \caption{\textbf{Cross-mouse consistency is not driven by shared
  adaptation, pupil-area fluctuations, or running-speed fluctuations.}
  \textbf{(A)} Per-image consistency in the chronologically early vs.\
  late trial half (scatter; dashed line: equal consistency).  Spearman
  $r = 0.960$ ($p = 6.1 \times 10^{-66}$) across images.  Mean early
  $= 0.275$; mean late $= 0.259$; mean difference early $-$ late
  $= +0.015$, in the direction opposite to a shared adaptation confound.
  \textbf{(B)} Cross-mouse pupil-area CV vs.\ per-image consistency
  (Spearman $r = +0.056$, $p = 0.55$, $n = 26$ sessions with
  eye-tracking).  No significant relationship.
  \textbf{(C)} Cross-mouse running-speed CV vs.\ per-image consistency
  (Spearman $r = -0.033$, $p = 0.72$, $n = 32$ sessions).  No
  significant relationship.}
  \label{fig:controls}
\end{figure}

\subsection{The transient is specific to natural scenes: gratings produce
a flat, still-rising response rather than a transient peak}

Static gratings produced substantially higher absolute consistency than
natural scenes (mean $r = 0.720$ vs.\ $r = 0.183$ for natural scenes in
the fixed 0--250~ms window; Figure~\ref{fig:gratings}), but the temporal
structure differed qualitatively.  Both stimulus classes have similarly
high within-mouse response reliability (gratings mean $= 0.992$, natural
scenes mean $= 0.978$; Methods), so this comparison is not confounded by
differential reliability between conditions.  Gratings rose by
$\sim$60~ms and then remained flat for the rest of the 0--250~ms epoch
rather than resolving into a clear plateau-then-hold: the quoted peak
($r = 0.680$) falls at the 250~ms window edge, where the curve is still
drifting slightly upward, giving a near-unity T/S ratio (approximately
$1.00\times$).  Natural scenes showed the characteristic
transient-plus-sustained structure (peak $r = 0.307$ at 100~ms,
T/S $= 1.68\times$).  We restrict quantitative comparisons to the
0--250~ms window for the same contamination reason as natural scenes.

The raw-magnitude gap above is not directly comparable, however: gratings
have 30 conditions (29-dimensional similarity profiles) while natural
scenes have 118 images (117-dimensional profiles), and the per-item
Spearman consistency measure's noise properties depend on this
dimensionality independently of any real representational difference. To
control for this, we recomputed natural-scenes consistency after
subsampling the same 30 images per random draw (50 draws, matching the
gratings condition count exactly; Figure~\ref{fig:gratings}, light blue).
This raised the fixed-window estimate from $r = 0.183$ to
$r = 0.303$ (95\% range across draws: $[0.258, 0.354]$) --- confirming
that part of the raw magnitude gap was a dimensionality artefact --- but
the item-count-matched value remains far below gratings ($r = 0.720$), and
the T/S ratio is essentially unchanged under matching ($1.66\times$ vs.\
$1.68\times$ unmatched).  The shape comparison between stimulus classes,
which is unaffected by item count, is therefore the more reliable
evidence for a genuine difference; the magnitude comparison should be
read as directionally consistent with it but partially inflated by the
dimensionality mismatch.

This comparison is consistent with the view that feedforward responses
to simple, low-complexity stimuli are sufficient to produce stable shared
representations in V1, while the complex stimulus-specific inhibitory
sharpening mechanism is only recruited for natural images that engage a
broader, less stereotyped neural population.

\begin{figure}[htbp]
  \centering
  \includegraphics[width=\textwidth]{figures/fig8_gratings_vs_natural_scenes.png}
  \caption{\textbf{Static gratings produce a weaker and less transient
  cross-mouse consistency than natural scenes.}
  Raw mean cross-mouse consistency over time for natural scenes (dark
  blue, peak $r = 0.307$ at 100~ms, T/S $= 1.68\times$) and static
  gratings (gold, rises by $\sim$60~ms then stays flat,
  T/S $\approx 1.00\times$).  Light blue: natural-scenes consistency
  recomputed after subsampling to 30 images per draw (50 draws, shaded
  band = 95\% range across draws), matching the gratings condition count
  to control for the item-count confound in the raw-magnitude comparison
  (peak $r = 0.291$, T/S $= 1.66\times$) --- closely tracking the
  full-118-image curve in shape while confirming that some of the raw
  magnitude gap reflects dimensionality rather than stimulus identity.
  Gratings rise to a near-immediate plateau; natural scenes show the
  transient-plus-sustained structure regardless of item count.  Dotted
  vertical line: stimulus offset (250~ms); data beyond this point risk
  contamination from the next stimulus.  Both conditions have similarly
  high within-mouse response reliability (gratings $= 0.992$, natural
  scenes $= 0.978$), so the shape difference is not a reliability
  artefact.}
  \label{fig:gratings}
\end{figure}

\subsection{Cross-area hierarchy: consistency and onset latency vary with
hierarchical position}

We applied the same sliding-window pipeline to four additional visual
areas (VISl, VISal, VISpm, VISam; session counts 24, 23, 20, 26
respectively, with $\geq 10$ QC-passing units per session).  VISp had the
highest peak consistency ($r = 0.307$) and earliest onset (50~ms, peak at
100~ms).  The four extrastriate areas showed later onsets and lower,
roughly comparable peak consistency rather than a monotonic decline with
hierarchical position: VISl $r = 0.186$ (onset 60~ms, peak 110~ms),
VISal $r = 0.234$ (onset 60~ms, peak 100~ms), VISpm $r = 0.210$
(onset 70~ms, peak 120~ms), VISam $r = 0.238$ (onset 70~ms, peak 110~ms;
Figure~\ref{fig:hierarchy}).  All five areas have similarly high
within-mouse response reliability (range 0.964--0.978; Methods), so
this ordering reflects genuine differences in raw consistency rather
than differential reliability.  The transient-plus-sustained temporal
structure was qualitatively present in all areas.  Between-area
comparisons are reported descriptively; areas share sessions, and a
formal test would require a mixed model accounting for this dependence.

The five areas also differ in mean per-session unit count (VISp 63.0,
VISl 38.9, VISal 67.1, VISpm 44.0, VISam 57.9), which inflates
similarity-matrix SNR and hence raw T/S ratio independently of any
genuine areal difference.  To control for this, we randomly subsampled
every area's unit pool down to 39 units/session (VISl's mean, the
smallest of the five) over 20 draws and recomputed T/S on each matched
draw.  Matched T/S ratios were uniformly higher than raw (e.g.\ VISp:
$1.68\times \to 1.81\times$; VISl: $2.03\times \to 2.13\times$; full
values in Methods), and the ranking was preserved for three of the five
areas (VISl $>$ VISam $>$ VISal, both raw and matched), but VISp and
VISpm swapped order: VISpm was marginally higher than VISp under the raw
comparison ($1.74\times$ vs.\ $1.68\times$) but marginally lower under
the matched comparison ($1.79\times$ vs.\ $1.81\times$).  This pair is
therefore not distinguishable once unit count is controlled, and the
overall cross-area ranking should be read as robust except for the
VISp/VISpm comparison specifically.

\begin{figure}[htbp]
  \centering
  \includegraphics[width=\textwidth]{figures/fig9_cross_area_hierarchy.png}
  \caption{\textbf{Cross-mouse representational consistency is highest in
  V1 and lower and roughly comparable across four extrastriate areas.}
  \textbf{(A)} Temporal profiles of cross-mouse consistency for all five
  visual areas (VISp: blue; VISl: green; VISal: yellow; VISpm: red;
  VISam: purple; $n$ per area in legend).  Stars: empirical peak per area.
  Transient-plus-sustained structure is present across the hierarchy;
  VISp shows the most pronounced transient.
  \textbf{(B)} Peak raw consistency (Spearman $r$) per area (bars, with
  T/S ratio annotated).  VISp shows the highest peak; the four
  extrastriate areas are lower and similar to one another rather than
  forming a monotonic gradient.  Note: areas share sessions; comparisons
  are descriptive.}
  \label{fig:hierarchy}
\end{figure}

% ── Discussion ────────────────────────────────────────────────────────────────
\section{Discussion}

We have shown that mouse V1 does not represent natural images with a
constant degree of cross-individual agreement.  Instead, it passes through
a transient, low-dimensional population state --- peaking at $\sim$100~ms
post-stimulus and visible directly in the participation ratio of the
shared representation --- and this transient state is precisely the moment
at which different animals' representations are most alike.  Three further
lines of evidence indicate that this transient compression is produced by
fast-spiking inhibitory interneurons acting through inhibition-dominated
feedforward dynamics.

\subsection{Summary of convergent evidence}

The two most direct findings require no cell-type classification, no
circuit model, and no auxiliary assumptions.  First, the temporal profile
of cross-individual consistency follows a stereotyped
transient-plus-sustained shape that is incompatible with purely feedforward
tuning (which would produce a step function) or purely recurrent attractor
dynamics (which would maintain a plateau).  Second, this transient peak
coincides precisely with a collapse of the shared representational geometry
to its minimum dimensionality (PR~$= 31.5$ at 100~ms, vs.\ a pre-stimulus
baseline of $\approx 108$, for the cross-mouse-averaged similarity
matrix), and PR and consistency are tightly coupled across the entire
post-stimulus epoch (Spearman $r = -0.853$; confirmed by a
non-overlapping-window control on independent windows, $r = -0.929$,
$p = 4.6\times10^{-6}$; Methods).  The same coupling holds,
essentially as strongly, when PR is computed independently within each
mouse and only then averaged (per-mouse PR: pre-stimulus $\approx 35.6 \to$
peak 15.3; Spearman $r = -0.805$), which rules out the concern that the
averaged-matrix result is merely a restatement of consistency itself.
These two
findings are measured directly from the same similarity matrices and are,
in effect, two views of a single underlying phenomenon: cross-individual
agreement is high exactly when the population code is compressed into a
small number of stereotyped dimensions. Because both quantities are also
plausibly driven by the shared stimulus-evoked response envelope, we
tested this coupling against a common-cause account by partialling out a
response-envelope proxy (mean raw population response per window); the
coupling was essentially unchanged (avg-matrix: partial
$r = -0.801$ vs.\ raw $-0.853$; per-mouse: partial $r = -0.718$ vs.\ raw
$-0.805$, both $p < 10^{-3}$), indicating that PR tracks consistency
beyond what shared response magnitude predicts on its own. Under an
autocorrelation-robust circular-shift test (Methods), the raw
envelope--consistency correlation was no longer distinguishable from
chance ($p = 0.211$), weakening rather than strengthening a common-cause
interpretation, while the PR--consistency coupling itself is confirmed
by a fully independent, non-overlapping-window control ($r = -0.929$,
$p = 4.6\times10^{-6}$). This coupling
is the strongest and most direct result in the paper, and it does not
depend on any of the mechanistic claims that follow.

The remaining results address \textit{why} this compression occurs.
Third, FSU units produce a sharper transient than RSU units despite
identical onset latency, a difference that survives a unit-count-matched
subsampling control (matched RSU T/S $= 2.22\times$, still below FSU's
$2.29\times$), and sessions with more FSU units show correspondingly
sharper population-level transients, a difference that likewise survives
matching the two session groups' unit counts (though roughly a quarter of
the raw absolute-consistency gap between them is attributable to a
similarity-matrix SNR advantage from unit count rather than a genuine FSU
effect).  Fourth, genetically confirmed PV$^+$ units show an even
sharper transient (T/S $= 3.59\times$) than waveform-inferred FSU units
(T/S $= 2.29\times$), strengthening the cell-type attribution while
revealing an impurity in the waveform classification.  Fifth, a
five-parameter Wilson-Cowan E/I model fit directly to the image-selective
firing-rate amplitude of training sessions generalises to held-out
sessions with mean test $R^2 = 0.877$ and recovers an inhibition-dominated
ratio ($r_\mathrm{EI} = 0.647 \pm 0.010$) stable across all five
cross-validation folds; an independent fit of the same model to
cross-mouse consistency curves converges on a similar ratio
($r_\mathrm{EI} = 0.608$), while a fit to uncorrected total firing rate
gives a markedly different, less stable estimate
($r_\mathrm{EI} = 0.924$), indicating the first two fits' agreement
reflects a real property of the data rather than curve-fitting
flexibility.

Among the mechanistic evidence, the cross-validated E/I model is the
strongest single piece, precisely because it does not rest on consistency
alone: fit directly to real, image-selective firing rates, a minimal
circuit description trained on $\sim$26 sessions reproduces the temporal
dynamics of the remaining $\sim$6 with test $R^2 > 0.80$ in every fold,
the inhibition-dominance conclusion is robust even in the worst fold, and
an independently-computed cross-mouse consistency curve is fit by the
same model to a similar ratio. The optotagging result, despite the
small sample ($n = 3$ usable sessions), is the most direct evidence for
the PV$^+$ attribution specifically.  We emphasise, however, that the
central empirical claim --- the existence and timing of the transient
low-dimensional state, and its coupling to cross-individual consistency
--- does not depend on the correctness of the circuit-level
interpretation.

\subsection{Relationship to prior work}

The high-dimensional spontaneous geometry of V1 (pre-stimulus per-mouse
PR~$\approx 35.6$) is consistent with \citet{stringer2019}, who
characterised the within-animal dimensionality of V1 population
responses as a high-dimensional power-law eigenspectrum; the per-mouse
PR, rather than the cross-mouse-averaged-matrix PR ($\approx 108$), is
the quantity comparable to their within-animal measurements, since
averaging similarity matrices across mice inflates apparent
dimensionality by washing out cross-mouse disagreement.  Our contribution
is the demonstration that stimulus onset drives a rapid collapse in both
the averaged-matrix PR (to $\approx 31.5$) and the per-mouse PR (to
$\approx 15.3$) --- a low-dimensional attractor by either measure --- and
that this collapse is shared across individuals.  The
subsequent partial recovery (PR $\to 61.2$ by 250~ms) suggests that the
sustained representation is less structured than the transient peak,
consistent with the smaller sustained consistency floor.

Our RSA approach differs from Procrustes alignment \citep{williams2021}
and hyperalignment \citep{haxby2011} in that it measures agreement on
representational \textit{structure} --- which stimuli are similar to which
others --- rather than the orientation of neural axes, making it invariant
to individual differences in which neurons are tuned to which features.
The mean consistency of $r = 0.314$, against a within-mouse RDM-space
reliability of $r = 0.832$ (Methods), indicates that a substantial share
of the available representational signal is shared across individuals,
with the remainder reflecting idiosyncratic tuning or developmental
variation.

The feedforward/recurrent distinction proposed by \citet{lamme2000} maps
naturally onto our two-component temporal profile: the transient at
$\sim$100~ms falls within the early feedforward epoch, while the sustained
floor likely reflects contributions from recurrent intra-areal and
inter-areal dynamics.  The fact that cross-mouse agreement is maximal
during the feedforward component is consistent with shared thalamic-driven
tuning being more stereotyped across individuals than recurrent
contributions shaped by individual experience and circuit variation.

\subsection{The optotagging complication and its interpretation}

The finding that only 6 of 18 optotagged PV$^+$ units fall within a
confidently-classified FSU waveform range, with 10 in an excluded
ambiguous band and 2 in the confidently-classified RSU range, is an
important caveat for waveform-based cell-type inference.  This result is
not unprecedented: prior work has noted that a subset of PV$^+$ cells
express intermediate waveform widths, particularly in superficial
cortical layers \citep{bartho2004,mccormick1985}.  Replacing a single
hand-picked 0.4~ms threshold with a two-component Gaussian mixture and a
symmetric per-component probability-mass bound (Methods) makes this
ambiguity explicit rather than hiding it: most borderline optotagged
units are now correctly excluded as ambiguous rather than being forced
into a bin, and only 2 of 18 are actively misclassified as RSU (versus 3
of 18 under the single-threshold rule).  The practical consequence for
our analysis is quantifiable: residual contamination of the RSU pool by
misclassified PV$^+$ cells inflates the RSU T/S ratio slightly and
deflates the FSU T/S ratio, producing exactly the pattern we observe
(optotagged T/S $3.59\times$ $>$ waveform FSU T/S $2.29\times$).  The
optotagging result therefore not only validates the direction of the
waveform-based finding but also provides a corrected, higher estimate of
the true PV$^+$ T/S ratio.  Our main conclusions are
\textit{strengthened} rather than weakened by this complication.

\subsection{Limitations}

\paragraph{E/I model identifiability.}
The Wilson-Cowan model is a phenomenological description, not a
mechanistic claim.  An earlier concern with this analysis was that the
model was fit only to cross-mouse consistency curves, an agreement
statistic rather than a firing rate, leaving open whether the recovered
$r_\mathrm{EI}$ reflected anything about actual population dynamics
rather than a five-parameter model's ability to match any pair of smooth
curves. We addressed this directly by fitting the same model to the
image-selective firing-rate amplitude of the RSU/FSU pools, computed
independently from the raw spike-count tensors; this fit converges on a
similar $r_\mathrm{EI}$ ($0.646$ vs.\ $0.608$ from consistency) and is, if
anything, better identified (SD $< 0.001$ across restarts vs.\ $0.120$
for the consistency fit; Supplementary Figure~S9C). A fit to uncorrected
total firing rate, which retains non-selective drive that the
selective-amplitude and consistency computations both exclude, gives a
markedly different and less stable estimate ($r_\mathrm{EI} = 0.924$),
showing that the model does discriminate between fit targets rather than
absorbing any input into the same answer. Some identifiability concerns remain specific to individual fits: the
$w_\mathrm{II}$ parameter converged to its upper bound in the consistency
fit, so its absolute value should not be interpreted there; $r_\mathrm{EI}$
itself is stable across widened $w_\mathrm{II}$ bounds for all three fit
targets (Supplementary Figure~S9D).  The 10-random-initialisation stability check
also reveals a competing, recurrence-dominated optimum
($r_\mathrm{EI} \approx 1.01$) that one of ten restarts converges to for
the consistency fit specifically; the inhibition-dominated solution is
the dominant basin there (9/10 restarts) but not the only one, so that
particular fit should be read as a strong majority result rather than a
globally unique optimum. More constrained inference would require
single-trial data or perturbational experiments that independently vary
E and I drive.

\paragraph{Sustained floor interpretation.}
The sustained component ($r = 0.182$, 220--250~ms) is estimated from a
narrow window forced by the back-to-back presentation protocol.  Whether
it reflects a genuine attractor state, residual stimulus-driven activity,
or early contributions from the subsequent image cannot be fully resolved.
A design with variable inter-stimulus intervals would be necessary to
distinguish these possibilities.

\paragraph{Optotagging sample size.}
The optotagged PV$^+$ analysis rests on only three sessions contributing
3, 3, and 10 confirmed units respectively.  The T/S estimate should be
treated as directional support for the waveform-based finding, not as a
precise quantitative estimate of the true PV$^+$ T/S ratio.

\paragraph{Residual unit-count confound.}
FSU, RSU, and the high-/low-FSU-fraction session groups all differ in
unit count, which inflates similarity-matrix SNR and hence raw
consistency independently of any genuine cell-type effect.
Unit-count-matched subsampling controls (Results) show that the FSU/RSU
T/S difference and the high-/low-FSU T/S difference both survive
matching, but they also show the confound is only partially, not fully,
resolved: the matched RSU T/S ($2.22\times$) sits closer to the FSU T/S
($2.29\times$) than the raw comparison suggests, and roughly a quarter of
the raw absolute-consistency gap between high- and low-FSU-fraction
sessions is attributable to unit count rather than FSU abundance itself.
The T/S metric's within-curve normalisation limits, but does not
eliminate, this confound, since unit count also affects the curve's
overall noise floor and hence, to a lesser degree, its shape.

\paragraph{Cross-area comparisons.}
Areas share sessions, so between-area comparisons of consistency levels
are not statistically independent.  The observed pattern --- highest in
VISp, lower and roughly comparable across the four extrastriate areas
--- is consistent with known hierarchical properties of the mouse visual
system but should be interpreted descriptively pending a formal
mixed-model analysis.  Areas also differ in mean per-session unit count,
which we addressed with a unit-count-matched subsampling control
(Results); the T/S ranking was preserved for three of the five areas but
VISp and VISpm swapped order between the raw and matched comparisons,
indicating that pair specifically is not distinguishable once unit count
is controlled.

\paragraph{Attractor geometry and image content.}
The two dominant attractor dimensions account for only 19.3\% of the
shared variance.  The remaining shared structure is not accounted for by
the low-level image statistics and CNN features examined here, and may
reflect higher-level image properties or individual-animal variation not
captured by our cross-mouse RSA measure.

\subsection{Broader significance}

These findings establish that the transient compression of the V1
population code is a measurable, reproducible property of the visual
response --- and, in our data, a mechanistically interpretable one.  The
cross-validated generalisability of the E/I model implies that the circuit
mechanism producing this trajectory is not merely a feature of the 32 mice
sampled but a stable property of the mouse V1 architecture.

More broadly, if the shared geometry of V1 is transiently imposed by
inhibitory interneurons during a 50--150~ms window after stimulus onset,
then downstream areas reading out V1 activity during this window receive
a cleaner, more stereotyped signal than at other times.  This temporal
gating of canonical representational structure may be a general principle
of cortical computation, extending beyond V1 to other sensory and
association areas where fast-spiking interneurons play similar roles in
shaping population dynamics.  The finding that the transient is absent for
static gratings and lower in extrastriate areas relative to VISp further
suggests that this mechanism is specifically engaged by complex,
naturalistic stimuli that require active representational compression,
though the drop from VISp to the four extrastriate areas does not itself
follow a graded, step-by-step pattern.

% ── Methods ───────────────────────────────────────────────────────────────────
\section{Methods}

\subsection{Dataset}

All data were obtained from the Allen Brain Observatory Visual Coding
Neuropixels dataset \citep{siegle2021}, accessed via the AllenSDK
(version~2.x).  We used all 32 sessions of type
\textit{brain\_observatory\_1.1}, in which head-fixed mice passively
viewed a standardised battery of natural and artificial stimuli while six
Neuropixels probes recorded simultaneously from multiple visual areas.
The natural scenes stimulus (118 greyscale images, 250~ms per image, 50
repetitions per image, back-to-back with no inter-stimulus interval) and
the static gratings stimulus (orientation $\times$ spatial-frequency
conditions, 250~ms per presentation) were used.  All analyses used the
VISp area designation from the Allen Common Coordinate Framework~v3 for
the main analysis; the cross-area analysis additionally included VISl,
VISal, VISpm, and VISam.  In the \textit{brain\_observatory\_1.1} dataset
there is one session per mouse, so ``32 mice'' and ``32 sessions'' are
used interchangeably throughout, and sessions are the resampling/held-out
unit in the cluster bootstrap and cross-validation procedures described
below.

\subsection{Unit quality control}

Units were retained if they met all three Allen Brain Observatory
recommended quality criteria: amplitude cutoff $< 0.1$, presence ratio
$> 0.9$, and ISI violation rate $< 0.5$.  Sessions were included in the
analysis of a given area only if they contributed at least 10 QC-passing
units in that area.  All 32 sessions contributed at least 10 VISp units
(range: 14--110).  Session counts for other areas were: VISl 24,
VISal 23, VISpm 20, VISam 26.  Mean units/session were: VISp 63.0,
VISl 38.9, VISal 67.1, VISpm 44.0, VISam 57.9.  To control for the
resulting SNR confound in cross-area T/S comparisons (Results), each
area's unit pool was randomly subsampled down to 39 units/session (VISl's
mean, the smallest of the five) over 20 draws, and the sliding
consistency curve and T/S ratio were recomputed on each matched draw.
Raw vs.\ matched T/S ratios were: VISp $1.68\times \to 1.81\times$
(95\% range across draws $1.73$--$1.89\times$), VISl
$2.03\times \to 2.13\times$ ($1.97$--$2.30\times$), VISal
$1.84\times \to 1.92\times$ ($1.84$--$2.00\times$), VISpm
$1.74\times \to 1.79\times$ ($1.72$--$1.87\times$), VISam
$2.00\times \to 2.10\times$ ($1.99$--$2.23\times$).

\subsection{Preprocessing and response matrix construction}

Spike times were binned into 1~ms bins spanning $-75$ to $+404$~ms
relative to stimulus onset and stored as a tensor of shape
(images $\times$ units $\times$ bins) per session.  Any response window
$[t_\mathrm{lo}, t_\mathrm{hi})$~ms was reconstructed by summing the
relevant bin columns, then z-scoring per unit (across images).  This
single-pass preprocessing avoids redundant NWB reads across analyses.
For the main 0--250~ms analysis, z-scored matrices were computed
directly from the summed tensor columns.

\subsection{Within-mouse reliability}

Within-mouse split-half reliability was estimated by randomly partitioning
each image's 50 trials into two halves, computing the mean population
response vector for each half, and measuring their Spearman rank
correlation.  This was repeated 20 times and averaged.  The Spearman-Brown
formula $r_\mathrm{SB} = 2r/(1+r)$ estimated full-data reliability.  This
response-vector reliability (mean of $r_\mathrm{SB}$ across images and
mice) is a QC diagnostic confirming the analysis is not power-limited by
within-session trial noise, but it is not in the same space as
cross-mouse consistency and should not be used to normalise it: it is
computed from raw per-unit response vectors, whereas consistency is
computed on image~$\times$~image similarity structure (RDM space).

To obtain a reliability estimate genuinely comparable to consistency, we
additionally computed an RDM-space split-half reliability for VISp
natural scenes, the one area/stimulus combination for which per-trial
spike data was available.  Each mouse's trials were split in half (20
random splits), each half was converted to a response matrix and then to
a cosine similarity matrix (identical to the main RSA pipeline below),
and the two resulting per-mouse similarity matrices were Spearman
rank-correlated against each other exactly as \texttt{per\_image\_consistency}
correlates two different mice's similarity matrices, then Spearman-Brown
corrected and averaged over splits and mice.  This RDM-space reliability
(mean $r = 0.832$) is the number reported alongside consistency values in
Results.  It was not computed for the four extrastriate areas or for
static gratings because per-trial spike data for those conditions was not
retained in the cached preprocessing archive; area and stimulus-type
comparisons therefore use raw consistency values throughout rather than a
normalised percentage, since both response-vector reliability (0.964--0.992
across all areas and stimulus types) and, where computable, RDM-space
reliability are high and broadly comparable across the conditions being
compared.

\subsection{Representational similarity analysis and cross-mouse
consistency}

For each session, the inter-image cosine similarity matrix was computed
from the z-scored response matrix:
$S_{ij} = \mathbf{x}_i \cdot \mathbf{x}_j /
(\|\mathbf{x}_i\| \|\mathbf{x}_j\| + \epsilon)$,
$\epsilon = 10^{-8}$.  Per-image cross-mouse consistency was the mean
Spearman rank correlation, across all mouse pairs $(a,b)$, between their
row-wise similarity profiles for each image.

Cluster bootstrap confidence intervals resampled whole mice (not pairs)
with replacement for 1000 iterations.  When a mouse appeared multiple
times in a bootstrap sample, pairs consisting of two copies of the same
original mouse were excluded, preventing spurious self-correlation from
inflating bootstrap consistency estimates.

\subsection{CNN embeddings}

ResNet50 pretrained on ImageNet (IMAGENET1K\_V2) was used as a feature
extractor; global average pooling was applied to each convolutional block
to yield a fixed-length embedding per image.  Ridge regression with
leave-one-out cross-validation (LOO, 118 images) was used because the
feature dimension (256--2048) greatly exceeds the number of images.
Layer~2 (512-dimensional) was selected as the primary CNN predictor on
the basis of LOO accuracy across layers (Figure~\ref{fig:predictors}C);
the final layer (1000-dimensional), which performed comparably, was
carried through the same sparsity-comparison analysis as a
robustness check (Figure~\ref{fig:predictors}D, Supplementary
Figure~\ref{fig:S2}).

Whether the CNN-based signal added explanatory power beyond sparsity
was assessed primarily by a nested $R^2$ comparison: $R^2$ for
consistency~$\sim$~sparsity alone was compared against
consistency~$\sim$~sparsity~$+$~CNN, where the CNN term was the
full-model ridge LOO prediction (Figure~\ref{fig:predictors}D). This
model-comparison statistic was treated as the headline evidence because,
unlike a partial correlation, it does not require residualising a
fitted, cross-validated prediction against the variable it was fit to
predict.

Two partialling methods were additionally computed as a secondary,
robustness check, reported in Supplementary Figure~\ref{fig:S2}.
(1)~\textit{Added-variable (OLS) partial correlation}: consistency,
sparsity, and the full-model CNN LOO prediction were each residualised
by ordinary least-squares regression on the other predictor, and the
Pearson correlation of the residuals gives the partial correlation in
each direction (sparsity~$|$~CNN and CNN~$|$~sparsity).
(2)~\textit{Rank-residualised (Spearman) partial correlation}, an
independent, non-parametric check reported in Supplementary
Figure~\ref{fig:S2}C: both variables were rank-transformed, the linear
effect of the rank-transformed covariate was regressed out, and the
Spearman correlation of the residuals was computed.  In both cases, the
CNN quantity partialled was the full ridge LOO prediction from the
512-dimensional layer~2 (or 1000-dimensional final-layer) embedding, not
a scalar summary of the embedding.  Because these partial-correlation
estimates residualise a prediction that was itself cross-validated
against consistency, and because their magnitude varied substantially
between the two partialling methods (Results), they are reported as
illustrative robustness detail rather than as the primary test of
independence.

\subsection{Sliding-window temporal analysis}

The RSA pipeline was repeated in 25~ms windows stepped at 10~ms
increments from $-50$ to $+370$~ms, reconstructed from the binned tensor.
Response matrices were z-scored independently within each window.  The
sustained floor was estimated from the 220--250~ms window only, to avoid
contamination from the subsequent image's V1 response.  The onset
threshold was defined as the first window where mean consistency exceeded
$\max(2 \times \text{pre-stimulus baseline},\, 0.05)$.

\subsection{Participation ratio}

The PR of a positive semi-definite matrix $\mathbf{M}$ with eigenvalues
$\lambda_i \geq 0$ is $\mathrm{PR}(\mathbf{M}) =
(\sum_i \lambda_i)^2 / \sum_i \lambda_i^2$.  At each time window,
the average similarity matrix was formed by summing per-session cosine
similarity matrices and dividing by the number of sessions.

\subsection{Response-envelope control}

Because both PR and cross-mouse consistency are computed from the same
per-window response matrices, we tested whether their covariation across
time is attributable to a shared response envelope rather than a specific
link between dimensionality and agreement. The envelope was defined as the
mean raw (pre-z-score) population response --- summed spike count per unit
per image within each window, averaged over units, images, and mice ---
computed upstream of the per-window z-scoring used for RSA, since that
normalisation removes each window's overall response magnitude by
construction. Given this envelope $z$, the relationship between consistency
$x$ and PR $y$ was re-assessed with the Spearman partial correlation
\begin{equation}
  r_{xy \cdot z} = \frac{r_{xy} - r_{xz} \, r_{yz}}
  {\sqrt{(1 - r_{xz}^2)(1 - r_{yz}^2)}},
  \label{eq:partial}
\end{equation}
with significance assessed via a $t$-test on $n - 3$ degrees of freedom.

\subsection{Autocorrelation-robust significance for sliding-window correlations}

Because the sliding windows overlap (25~ms windows stepped by 10~ms, or
20~ms as subsampled for the PR analysis), correlations computed across
post-stimulus windows do not have $n$ independent observations, and the
analytic $p$-values reported alongside them understate the true
$p$-value. We assessed the PR--consistency and envelope correlations
with two additional, assumption-appropriate tests.

First, a circular-shift permutation test: the null distribution was built
by circularly shifting one series relative to the other by every one of
the $n-1$ non-trivial lags, which preserves each series' own
autocorrelation structure while destroying only their mutual alignment;
the $p$-value is the fraction of shifted correlations at least as extreme
as the observed one. For the partial correlations, consistency was
shifted relative to the paired (PR, envelope), so the PR--envelope
relationship was preserved under the null. With $n = 19$ post-stimulus
windows this test has only 18 possible shifts, giving a resolution floor
of $p \approx 1/19 = 0.053$; several of our correlations were more
extreme than all 18 shifts and are reported at this floor rather than as
non-significant.

Second, a non-overlapping-window control: PR and consistency were
recomputed at windows spaced one full window-width (30~ms) apart, so no
two windows share any spike-count bins and the resulting points are
independent by construction. This yields fewer points (13 post-stimulus
windows) but permits a conventional analytic significance test with no
autocorrelation correction.

\subsection{Waveform-based cell class classification}

Units were classified as FSU (fast-spiking, putative PV$^+$) or RSU
(regular-spiking, putative pyramidal/excitatory) using trough-to-peak
waveform duration, following the canonical narrow/broad distinction
\citep{bartho2004,mccormick1985}, but via a two-threshold rule rather
than the single hand-picked cutoff used in that convention.  A
2-component Gaussian mixture model was fit to log waveform duration
across all 3,694 QC-passing VISp units (log-space fitting prevents a
small tail of long-duration, likely triphasic/artifact units from being
fit as a spurious third mode).  A unit was assigned to a class only if
(i)~its posterior probability of membership in the corresponding
component exceeded 0.9, and (ii)~it fell within a symmetric interval
capturing 60\% of that component's own probability mass around its mean
(0.84 standard deviations in log-space); units failing either criterion
were labelled ambiguous and excluded from both classes.  This yielded an
FSU range of 0.218--0.316~ms and an RSU range of 0.549--0.728~ms
(28.2\% of units excluded as ambiguous).  Sessions with fewer than 8
QC-passing units of a given class were excluded from that class's
analysis (FSU: 15/32 sessions retained; RSU: 31/32).  To control for the
resulting unit-count difference between classes, RSU sessions were also
subsampled down to the mean FSU unit count (8/session) over 50 random
draws, and the sliding consistency curve and T/S ratio were recomputed
on each matched draw (Results).

\subsection{Optotagging validation}

Genetically confirmed PV$^+$ units were identified from the five
Pvalb-IRES-Cre/wt$\times$Ai32(RCL-ChR2(H134R)\_EYFP)/wt sessions in
the \textit{brain\_observatory\_1.1} dataset.  Single light pulse epochs
were used.  QC-passing VISp units were classified as optotagged PV$^+$
if: (i)~mean spike count in the 0--5~ms post-pulse window $> 0.05$ spikes
per pulse; and (ii)~light-to-baseline ratio $> 2.0$, where baseline was
the mean spike count in the 5~ms pre-pulse window.  The 5~ms response
window excludes polysynaptic responses (pyramidal cells driven by PV$^+$
firing respond at 7--15~ms latency).  Sessions with fewer than 3
optotagged VISp units were excluded from the cross-mouse RSA (sessions
746083955 and 773418906 contributed one unit each).

\subsection{E/I model}

The two-population Wilson-Cowan model \citep{wilson1972} is described by
equations~\eqref{eq:E}--\eqref{eq:I}.  The thalamic drive $D(t)$ is a
gamma-shaped function (shape parameter $k = 5$, scale $\theta = 14$~ms)
with onset at 40~ms, normalised to unit amplitude.  Time constants were
fixed at $\tau_E = 20$~ms and $\tau_I = 10$~ms.  The five free
parameters are $\{r_\mathrm{EI}, w_\mathrm{IE}, w_\mathrm{EI},
w_\mathrm{II}, I_\mathrm{amp}\}$ with bounds $r_\mathrm{EI} \in
[0.01, 2.0]$, $w_\mathrm{IE} \in [0.1, 20.0]$,
$w_\mathrm{EI} \in [0.1, 15.0]$, $w_\mathrm{II} \in [0.0, 8.0]$,
$I_\mathrm{amp} \in [0.5, 20.0]$.  The model was integrated with
\texttt{scipy.integrate.odeint} (rtol $= 10^{-6}$, atol $= 10^{-8}$)
and fit by minimising the sum of mean squared errors between
normalised model outputs and normalised RSU/FSU curves simultaneously
using L-BFGS-B, with ten initialisations (one informed, nine random) per
fit target.  The model was fit independently to three targets: (i)
cross-mouse consistency curves, (ii) uncorrected total firing rate ---
the mean per-unit, per-image spike rate (Hz) in each sliding window,
summed across the pool --- and (iii) image-selective firing-rate
amplitude, the per-unit standard deviation of firing rate across images
(Hz), averaged over the pool, in the same sliding windows.  Rate curves
(ii) and (iii) were computed directly from the raw spike-count tensors
using the same RSU/FSU unit classification, sessions, and window grid as
the consistency curves.  Sensitivity to the $w_\mathrm{II}$ upper bound,
shown separately for all three fit targets, is in Supplementary
Figure~S9D.

\subsection{Cross-validation}

Five-fold cross-validation was performed over sessions, for the
selective-amplitude fit.  In each fold, RSU and FSU selective-amplitude
curves were computed from the $\sim$80\% training sessions; the model was
fit using three warm starts near the full-data solution.  The fitted
parameters were then used to predict curves from the $\sim$20\% held-out
sessions.  Train and test $R^2$ were computed on amplitude-normalised
curves.

\subsection{Grating response matrices}

Grating conditions were defined as unique (orientation, spatial frequency)
pairs; only conditions present in all 32 sessions (intersection:
30~conditions) were retained.  The same QC thresholds, z-scoring
procedure, and RSA pipeline were applied as for natural scenes.

\subsection{Statistical tests}

Spearman rank correlation was used throughout for robustness to outlier
images and bounded values.  Wilcoxon signed-rank tests were used for
paired comparisons of bounded variables.  No multiple comparisons
correction was applied to the two pre-specified behavioural state tests
(pupil CV and running CV), which are independent by construction.  All
bootstrap confidence intervals used 1000 iterations.  The global random
seed was set to 42 and passed explicitly to all stochastic operations.

\subsection{Data and code availability}

All neural data are publicly available through the Allen Brain Observatory
(\url{https://portal.brain-map.org}) and accessed using the AllenSDK
Python package.  Analysis code is available at
\url{https://github.com/Jana-HQ/cross-mouse-v1-rsa}.

% ── Acknowledgements ──────────────────────────────────────────────────────────
\section*{Acknowledgements}

The author thanks the Allen Institute for Brain Science for making the
Visual Coding Neuropixels dataset publicly available.  This research did
not receive any specific grant from funding agencies in the public,
commercial, or not-for-profit sectors.

% ── Competing interests ───────────────────────────────────────────────────────
\section*{Competing interests}

The author declares no competing interests.

% ── Supplementary figures ──────────────────────────────────────────────────────
\newpage
\section*{Supplementary figures}

\begin{figure}[htbp]
  \centering
  \includegraphics[width=\textwidth]{figures/figS1_image_gallery.png}
  \caption{\textbf{Gallery of high- and low-consistency natural images.}
  Top and bottom 8 images ranked by cross-mouse consistency.  Consistency
  values shown in subtitles.  High-consistency images tend to be sparse and
  contain a single dominant feature; low-consistency images are denser and
  more spatially heterogeneous.  This is a qualitative illustration; no
  statistical test is performed on image content.}
  \label{fig:S1}
\end{figure}

\begin{figure}[htbp]
  \centering
  \includegraphics[width=\textwidth]{figures/figS2_cnn_detail.png}
  \caption{\textbf{Partial-correlation robustness checks for the
  CNN--sparsity relationship, layer~2 and the final ResNet50 layer ---
  see Figure~\ref{fig:predictors}D for the headline nested-$R^2$
  result.}  Layer~2 (purple circles) and the final layer (teal circles)
  are shown together in every panel to confirm the pattern is not
  specific to one layer.
  \textbf{(A)} Added-variable residual scatter, sparsity~$|$~CNN
  (layer~2: partial $r = -0.523$; final: partial $r = -0.424$).
  \textbf{(B)} Added-variable residual scatter, CNN~$|$~sparsity
  (layer~2: partial $r = -0.501$; final: partial $r = -0.378$).
  \textbf{(C)} Rank-based (Spearman) partial correlation, CNN~$|$~sparsity
  (layer~2: $r = -0.828$, $p = 5.5\times10^{-31}$; final: $r = -0.887$,
  $p = 9.9\times10^{-41}$) --- a substantially larger-magnitude partial
  correlation than the added-variable estimates in (A)--(B) and than the
  raw LOO $r$, a signature of sparsity acting as a partial suppressor: it
  shares variance with both consistency and the CNN prediction, so
  removing it from both can inflate residual alignment.  Directionally
  consistent with (A)--(B), but the magnitude swing across methods means
  this value should be read as illustrative rather than as a precise
  effect size (Methods; Figure~\ref{fig:predictors}D reports the robust
  nested-$R^2$ estimate).
  \textbf{(D)} Nested-$R^2$/partial-$r$ summary table comparing the
  added-variable (OLS) and rank-residualisation (Spearman) methods side
  by side, for both layers.}
  \label{fig:S2}
\end{figure}

\begin{figure}[htbp]
  \centering
  \includegraphics[width=\textwidth]{figures/figS3_early_vs_late.png}
  \caption{\textbf{Early vs.\ late window consistency (coarse robustness
  check).}
  Per-image consistency in the early feedforward epoch (0--75~ms) vs.\ the
  late epoch (75--250~ms; dashed line: equal consistency).  Early mean
  $= 0.084$, late mean $= 0.320$; Wilcoxon signed-rank
  $p = 4.2 \times 10^{-21}$; image rankings preserved (Spearman
  $r = 0.365$).  This is a coarser, non-overlapping-window version of
  the claim made by the full temporal profile in Figure~\ref{fig:temporal};
  it is included as an independent robustness check.}
  \label{fig:S3}
\end{figure}

\begin{figure}[htbp]
  \centering
  \includegraphics[width=\textwidth]{figures/figS4_pv_waveform_confusion.png}
  \caption{\textbf{Waveform-based FSU classification under-detects
  genetically confirmed PV$^+$ interneurons.}
  \textbf{(A)} Strip plot of optotagged PV$^+$ unit waveform durations by
  PV-Cre session ($n = 5$ sessions, 18 optotagged units total).  Shaded
  band and dashed lines: GMM narrow/broad boundaries (0.316~ms and
  0.549~ms).  Two units fall above the broad boundary (confidently-RSU);
  ten fall in the excluded ambiguous band.
  \textbf{(B)} Confusion matrix: optotagging outcome vs.\ GMM-inferred
  class --- FSU, ambiguous, RSU --- (PV-Cre sessions only).  Sensitivity
  (PV$^+$ $\to$ GMM FSU) $= 33\%$; precision (GMM FSU $\to$ PV$^+$)
  $= 19\%$; most optotagged units (10/18) fall in the excluded ambiguous
  band rather than being confidently but wrongly classified.}
  \label{fig:S4}
\end{figure}

\begin{figure}[htbp]
  \centering
  \includegraphics[width=\textwidth]{figures/figS5_eigenvec_properties.png}
  \caption{\textbf{Eigenvector loadings on image properties (full
  heatmap).}
  Spearman $r$ between eigenvector projections and image properties
  (sparsity, magnitude, consistency, CNN layers~1--4) for the top 8
  dimensions of the average cosine similarity matrix at 100~ms.
  Asterisks: $p < 0.05$.  Dim~1 reflects signal strength and early CNN
  features; Dim~2 reflects sparsity and spatial scale.}
  \label{fig:S5}
\end{figure}

\begin{figure}[htbp]
  \centering
  \includegraphics[width=\textwidth]{figures/figS6_attractor_gallery.png}
  \caption{\textbf{Attractor-space image gallery.}
  Images at the extremes of Dim~1 (signal-strength axis) and Dim~2
  (sparsity axis) of the average similarity matrix at 100~ms (top 6 and
  bottom 6 projections per dimension).  Dim~1 high: high-magnitude,
  low-spatial-frequency images.  Dim~2 low: fine-grained,
  high-frequency images.}
  \label{fig:S6}
\end{figure}

\begin{figure}[htbp]
  \centering
  \includegraphics[width=\textwidth]{figures/figS7_rank_stability.png}
  \caption{\textbf{Image-ranking stability across the response epoch.}
  \textbf{(A)} Spearman rank-correlation matrix of per-image consistency
  across six epochs ($-25$, 25, 50, 100, 175, 250~ms).  Image rankings are
  near-zero before 50~ms, correlated between peak and sustained epochs
  (100 vs.\ 250~ms: $r = 0.652$), and near-zero between pre-stimulus and
  any post-stimulus epoch.
  \textbf{(B)} Rolling Spearman correlation of per-image consistency at
  each sliding window vs.\ the 100~ms reference.  Significance coding
  (filled marker: $p < 0.05$, hollow marker: n.s.) uses uncorrected
  $p$-values for illustration; 43 overlapping windows are not independent.
  See Methods for an autocorrelation-robust reassessment of an analogous
  overlapping-window correlation.}
  \label{fig:S7}
\end{figure}

\begin{figure}[htbp]
  \centering
  \includegraphics[width=\textwidth]{figures/figS8_second_order_rsa.png}
  \caption{\textbf{Second-order RSA recovers the same temporal peak, but
  is a related rather than an independent statistic.}
  \textbf{(A)} Distribution of second-order Spearman $r$ across 118 images
  (mean $= 0.187$; all images $> 0$).  Second-order $r$ measures agreement
  between mice on how each image's full representational profile --- its
  similarity to all other images across five time windows --- changes over
  time.
  \textbf{(B)} Per-window mean second-order $r$ (bars $\pm$ SEM across
  images).  The peak window (100~ms) yields the highest value
  ($r = 0.307$), matching the first-order consistency peak exactly.
  \textbf{(C)} Second-order $r$ vs.\ first-order consistency
  (Spearman $r = 0.950$, $p = 2.4 \times 10^{-60}$): the near-redundancy
  of the two measures indicates second-order RSA is a re-parameterisation
  that tracks the same underlying signal as first-order consistency,
  rather than an orthogonal confirmation of it.}
  \label{fig:S8}
\end{figure}

\begin{figure}[htbp]
  \centering
  \includegraphics[width=\textwidth]{figures/figS9.png}
  \caption{\textbf{E/I model fit-target comparison and robustness checks.}
  \textbf{(A)} Model fit (dashed) vs.\ data (solid) for RSU
  ($R^2 = 0.901$) and FSU ($R^2 = 0.913$) consistency curves.
  $r_\mathrm{EI} = 0.608$.
  \textbf{(B)} Model fit vs.\ data for RSU ($R^2 = 0.805$) and FSU
  ($R^2 = 0.601$) uncorrected total firing rate. $r_\mathrm{EI} = 0.924$.
  \textbf{(C)} $r_\mathrm{EI}$ at convergence across 10 random
  initialisations, for each of the three fit targets. The
  selective-amplitude fit is essentially uniquely identified
  (SD $< 0.001$); the consistency fit is bimodal (9/10 restarts converge
  to $r_\mathrm{EI} \approx 0.608$, one to a competing basin near
  $r_\mathrm{EI} \approx 1.01$; mean $= 0.648$, SD $= 0.120$); the
  total-rate fit clusters near $r_\mathrm{EI} \approx 0.92$--$0.93$
  (mean $= 0.926$, SD $= 0.006$).
  \textbf{(D)} Sensitivity of fitted $r_\mathrm{EI}$ to the upper bound
  on $w_\mathrm{II}$ (bounds tested: 6, 8, 10, 12, 15), shown separately
  for all three fit targets: $r_\mathrm{EI}$ is stable across this range
  for each target, confirming the inhibition-dominance conclusion (where
  applicable) is not an artefact of the $w_\mathrm{II}$ constraint.}
  \label{fig:S9}
\end{figure}

\begin{figure}[htbp]
  \centering
  \includegraphics[width=\textwidth]{figures/figS10_adaptation_control.png}
  \caption{\textbf{Adaptation control extended.}
  \textbf{(A)} Distribution of per-image (early $-$ late) consistency
  differences (mean $= +0.015$; one-sample $t$-test $p = 1.1 \times
  10^{-8}$).  Positive mean indicates slight attenuation in the second
  half, opposite to a shared adaptation confound.
  \textbf{(B)} Early-half per-image consistency vs.\ full-session
  consistency ($r = 0.985$, $p = 2.6 \times 10^{-90}$).
  \textbf{(C)} Late-half per-image consistency vs.\ full-session
  consistency ($r = 0.984$, $p = 2.4 \times 10^{-88}$).  Image rankings
  are stable across halves; the small absolute attenuation is consistent
  with non-selective neural fatigue.}
  \label{fig:S10}
\end{figure}

% ── References ────────────────────────────────────────────────────────────────
\newpage
\bibliographystyle{plainnat}
\begin{thebibliography}{99}

\bibitem[Barth{\'o} et al.(2004)]{bartho2004}
Barth{\'o}, P., Hirase, H., Monconduit, L., Zugaro, M., Harris, K.~D.,
and Buzs{\'a}ki, G. (2004).
Characterization of neocortical principal cells and interneurons by
network interactions and extracellular features.
\textit{Journal of Neurophysiology}, 92(1):600--608.

\bibitem[Cadena et al.(2019)]{cadena2019}
Cadena, S.~A., Denfield, G.~H., Walker, E.~Y., Gatys, L.~A.,
Tolias, A.~S., Bethge, M., and Ecker, A.~S. (2019).
Deep convolutional models improve predictions of macaque V1 responses to
natural images.
\textit{PLOS Computational Biology}, 15(4):e1006897.

\bibitem[Freund and Buzs{\'a}ki(1996)]{freund1996}
Freund, T.~F. and Buzs{\'a}ki, G. (1996).
Interneurons of the hippocampus.
\textit{Hippocampus}, 6(4):347--470.

\bibitem[Girard et al.(2001)]{girard2001}
Girard, P., Hup{\'e}, J.~M., and Bullier, J. (2001).
Feedforward and feedback connections between areas V1 and V2 of the
monkey have similar rapid conduction velocities.
\textit{Journal of Neurophysiology}, 85(3):1328--1331.

\bibitem[Haxby et al.(2011)]{haxby2011}
Haxby, J.~V., Guntupalli, J.~S., Connolly, A.~C., Halchenko, Y.~O.,
Conroy, B.~R., Gobbini, M.~I., Hanke, M., and Ramadge, P.~J. (2011).
A common, high-dimensional model of the representational space in human
ventral temporal cortex.
\textit{Neuron}, 72(2):404--416.

\bibitem[Kriegeskorte et al.(2008)]{kriegeskorte2008}
Kriegeskorte, N., Mur, M., and Bandettini, P.~A. (2008).
Representational similarity analysis -- connecting the branches of
systems neuroscience.
\textit{Frontiers in Systems Neuroscience}, 2:4.

\bibitem[Lamme and Roelfsema(2000)]{lamme2000}
Lamme, V.~A.~F. and Roelfsema, P.~R. (2000).
The distinct modes of vision offered by feedforward and recurrent
processing.
\textit{Trends in Neurosciences}, 23(11):571--579.

\bibitem[McCormick et al.(1985)]{mccormick1985}
McCormick, D.~A., Connors, B.~W., Lighthall, J.~W., and Prince, D.~A.
(1985).
Comparative electrophysiology of pyramidal and sparsely spiny stellate
neurons of the neocortex.
\textit{Journal of Neurophysiology}, 54(4):782--806.

\bibitem[Siegle et al.(2021)]{siegle2021}
Siegle, J.~H., Jia, X., Durand, S., Gale, S., Bennett, C.,
Graddis, N., Heller, G., Ramirez, T.~K., Choi, H., Luviano, J.~A.,
Groblewski, P.~A., et al. (2021).
Survey of spiking in the mouse visual system reveals functional hierarchy.
\textit{Nature}, 592(7852):86--92.

\bibitem[Somers et al.(1995)]{somers1995}
Somers, D.~C., Nelson, S.~B., and Sur, M. (1995).
An emergent model of orientation selectivity in cat visual cortical simple
cells.
\textit{Journal of Neuroscience}, 15(8):5448--5465.

\bibitem[Stringer et al.(2019)]{stringer2019}
Stringer, C., Pachitariu, M., Steinmetz, N., Carandini, M., and
Harris, K.~D. (2019).
High-dimensional geometry of population responses in visual cortex.
\textit{Nature}, 571(7765):361--365.

\bibitem[Swadlow(2003)]{swadlow2003}
Swadlow, H.~A. (2003).
Fast-spike interneurons and feedforward inhibition in awake sensory
neocortex.
\textit{Cerebral Cortex}, 13(1):25--32.

\bibitem[Troyer et al.(1998)]{troyer1998}
Troyer, T.~W., Krukowski, A.~E., Priebe, N.~J., and Miller, K.~D. (1998).
Contrast-invariant orientation tuning in cat visual cortex: feedforward
tuning and correlation-based intracortical connectivity.
\textit{Journal of Neuroscience}, 18(15):5908--5927.

\bibitem[Williams et al.(2021)]{williams2021}
Williams, A.~H., Kunz, E., Kornblith, S., and Linderman, S.~W. (2021).
Generalized shape metrics on neural representations.
\textit{Advances in Neural Information Processing Systems},
34:4738--4750.

\bibitem[Wilson and Cowan(1972)]{wilson1972}
Wilson, H.~R. and Cowan, J.~D. (1972).
Excitatory and inhibitory interactions in localized populations of model
neurons.
\textit{Biophysical Journal}, 12(1):1--24.

\bibitem[Yamins and DiCarlo(2016)]{yamins2016}
Yamins, D.~L.~K. and DiCarlo, J.~J. (2016).
Using goal-driven deep learning models to understand sensory cortex.
\textit{Nature Neuroscience}, 19(3):356--365.

\end{thebibliography}

\end{document}